% Book 4, Year 2, chapter 2 -- Copying, moving and smart pointers.
% Provenance (never printed): sources/book4/BRIEFS.md, chapter 2.
\chapter{Copying, Moving and Smart Pointers}\label{ch:b2:cpp-copy-move}

A function builds a vector of a million readings and returns it. In C the question would not arise: a function
returns a pointer and the caller frees the block. In C++ the vector is an object returned by value, and an object
returned by value is, in principle, a copy. The program below counts. Its element type \code{Tracked} (shown in
\cref{sec:b2:cpp-copy-move:moving}) adds one to a global counter each time one of its values is copied or moved.

\omcode{code/bachelor-2/02-cpp-copy-move/listings/hook.cpp}{6}{26}{A million values returned by
value}\label{lst:b2:cpp-copy-move:hook}

\omoutput{code/bachelor-2/02-cpp-copy-move/outputs/hook.txt}{Returning copied nothing; asking for a copy copied a
million values.}\label{out:b2:cpp-copy-move:hook}

Returning the vector copied nothing and moved nothing: the elements the caller holds are the very ones the
function built, at the same address. Writing \code{backup = r} copied all million. This chapter explains both
lines: what a copy of an object that owns memory must do, how C++ hands over the memory instead of copying it when
the source is about to die, and how the standard \omterm{def:b2:cpp-copy-move:unique-ptr}{smart pointers} express who owns a heap object.

\begin{recall}
\cref{ch:b2:cpp-classes-raii}: an object that owns a resource acquires it in its \omterm{def:b2:cpp-classes-raii:constructor}{constructor} and releases it in its
\omterm{def:b2:cpp-classes-raii:constructor}{destructor} (RAII), and copying such an object was forbidden with \code{= delete} until this chapter. A dynamic array
(Book 3) keeps $n$ elements in a block of capacity $c \geq n$; when it is full it allocates a larger block, carries
the $n$ elements over and frees the old one. Doubling the capacity makes the amortised cost of an append constant.
\end{recall}

\section{Copying an object that owns memory}\label{sec:b2:cpp-copy-move:copying}

\begin{definition}[Copy constructor, copy assignment operator]\label{def:b2:cpp-copy-move:copy}
The \emph{copy constructor}\index{copy constructor} of a class \code{T} is the \omterm{def:b2:cpp-classes-raii:constructor}{constructor} \code{T(const T &)}: it
builds a new object from an existing one, as in \code{T b = a;} or when \code{a} is passed by value. The \emph{copy
assignment operator}\index{copy assignment operator} \code{T &operator=(const T &)} gives an existing object the value
of another, as in \code{b = a;}. If a class declares neither, the compiler writes them: they copy each member in turn.
\end{definition}

For a class whose members are numbers and strings, the compiler's copy is right. For a class that holds a raw
pointer to a block it owns, it copies the pointer and not the block. The \code{Buffer} below owns a heap copy of a
street name; it has a \omterm{def:b2:cpp-classes-raii:constructor}{destructor} and lets the compiler write its copy.

\omcode{code/bachelor-2/02-cpp-copy-move/listings/shallow.cpp}{6}{26}{A buffer with the compiler's
copy}\label{lst:b2:cpp-copy-move:shallow}

\omoutput[1-6,11-15,18-21,24-26]{code/bachelor-2/02-cpp-copy-move/outputs/shallow.txt}{AddressSanitizer: the same
block freed twice.}\label{out:b2:cpp-copy-move:shallow}

The line printed by \code{main} is missing: the output stream had not been flushed when the sanitizer stopped the
program. The report itself is exact: the block allocated on line 22 (for \code{a}) was freed by a \omterm{def:b2:cpp-classes-raii:constructor}{destructor} called
on line 26 (that of \code{b}, destroyed first) and then again by the next \omterm{def:b2:cpp-classes-raii:constructor}{destructor} (that of \code{a}). Without the
sanitizer, freeing a block twice is undefined behaviour, which in practice corrupts the allocator's lists and fails
far from the cause.

\begin{definition}[Shallow copy, deep copy]\label{def:b2:cpp-copy-move:deep}
A \emph{shallow copy}\index{shallow copy} of an object copies its pointers, so that the copy and the original share
the blocks they point to. A \emph{deep copy}\index{deep copy} allocates new blocks for the copy and copies their
contents, so that the two objects share nothing and each can change or free its blocks alone.
\end{definition}

\begin{omfigure}
\begin{tikzpicture}[font=\footnotesize]
\node[font=\footnotesize\bfseries, anchor=west] at (-0.4,2.1) {shallow copy};
\foreach \o/\y in {a/1.2, b/0} {
  \node[cell, minimum width=1.1cm] (\o n) at (0.6,\y) {12};
  \node[cell, minimum width=1.1cm, right=0mm of \o n] (\o p) {$\bullet$};
  \node[anchor=east] at (0,\y) {\texttt{\o}};
}
\node[draw, fill=omThm!10, minimum width=2.8cm, minimum height=7mm, font=\ttfamily\scriptsize] (blk) at (4.6,0.6) {Harbour Road\textbackslash 0};
\draw[ptr] (ap.center) -- (blk.west);
\draw[ptr] (bp.center) -- (blk.west);
\node[font=\scriptsize, text=omThm, align=left, anchor=north] at (4.6,0.1) {one block, two owners:\\ freed twice};
\begin{scope}[xshift=7.6cm]
\node[font=\footnotesize\bfseries, anchor=west] at (-0.4,2.1) {deep copy};
\foreach \o/\y in {a/1.2, b/0} {
  \node[cell, minimum width=1.1cm] (\o n) at (0.6,\y) {12};
  \node[cell, minimum width=1.1cm, right=0mm of \o n] (\o p) {$\bullet$};
  \node[anchor=east] at (0,\y) {\texttt{\o}};
  \node[draw, fill=omMeth!10, minimum width=2.8cm, minimum height=7mm, font=\ttfamily\scriptsize] (blk\o) at (4.3,\y) {Harbour Road\textbackslash 0};
  \draw[ptr] (\o p.center) -- (blk\o.west);
}
\end{scope}
\end{tikzpicture}
\par\smallskip{\small Two \code{Buffer} objects after \code{Buffer b = a;}. Each object holds the length of the
text, 12, and a pointer to a block of 13 bytes. On the left, the compiler's copy: both pointers name the
same block. On the right, a \omterm{def:b2:cpp-copy-move:copy}{copy constructor} that allocates: each object owns its block.}
\end{omfigure}

The fix writes the \omterm{def:b2:cpp-copy-move:copy}{copy constructor} and the copy assignment by hand. The \omterm{def:b2:cpp-classes-raii:constructor}{constructor} allocates a block of its own
and copies the characters; the assignment is written in the form of \cref{sec:b2:cpp-copy-move:value}.

\omcode{code/bachelor-2/02-cpp-copy-move/listings/buffer.hpp}{7}{32}{A buffer with a \omterm{def:b2:cpp-copy-move:deep}{deep
copy}}\label{lst:b2:cpp-copy-move:buffer}

\omoutput{code/bachelor-2/02-cpp-copy-move/outputs/deep_copy.txt}{Three buffers, three blocks; the sanitizer is
silent.}\label{out:b2:cpp-copy-move:deep-copy}

Changing \code{b} left \code{a} alone, and the sanitizer found nothing to report. The three functions written here,
the \omterm{def:b2:cpp-classes-raii:constructor}{destructor}, the \omterm{def:b2:cpp-copy-move:copy}{copy constructor} and the copy assignment, go together: a class that needs one of them almost
always needs the three, since a \omterm{def:b2:cpp-classes-raii:constructor}{destructor} that frees means a raw resource that a member-wise copy would share.
This is the \emph{rule of three}, which \cref{sec:b2:cpp-copy-move:rules} completes.

\section{Value semantics}\label{sec:b2:cpp-copy-move:value}

\begin{definition}[Value semantics, reference semantics]\label{def:b2:cpp-copy-move:value-semantics}
A type has \emph{value semantics}\index{value semantics} if a copy of an object is a new, independent object with
the same value: changing one never changes the other, as with \code{int}, \code{std::string} and the
\code{Buffer} of \cref{lst:b2:cpp-copy-move:buffer}. A type has \emph{reference semantics}\index{reference
semantics} if copying an object gives a second name for the same underlying data, as with a pointer, or with every
object in Python.
\end{definition}

The C++ standard library gives its containers \omterm{def:b2:cpp-copy-move:value-semantics}{value semantics}, and a class that owns memory should do the same
unless there is a reason not to. The assignment of \cref{lst:b2:cpp-copy-move:buffer} has a pattern of its own,
named after its two steps.

\begin{method}[Copy and swap]\label{met:b2:cpp-copy-move:copy-and-swap}
To write \code{T &operator=(const T &o)} for a class that owns resources:
\begin{enumerate}
\item make a local copy \code{tmp} of \code{o} with the \omterm{def:b2:cpp-copy-move:copy}{copy constructor} (all the work that may fail happens here,
before \code{*this} is touched);
\item exchange the members of \code{*this} and \code{tmp}, with operations that cannot throw (\code{std::swap} of
pointers and numbers);
\item return \code{*this}; at the end of the function, \code{tmp} is destroyed and frees the old resources.
\end{enumerate}
\end{method}

\begin{proposition}[Copy and swap is strong]\label{prop:b2:cpp-copy-move:copy-and-swap}
An assignment written by \cref{met:b2:cpp-copy-move:copy-and-swap} gives the strong guarantee: if it throws, the
target is unchanged. It is also correct when the source and the target are the same object.
\end{proposition}

\begin{proof}
Only step 1 can throw, since swaps of pointers and integers cannot. If it throws, the \omterm{def:b2:cpp-copy-move:copy}{copy constructor} has cleaned
up after itself (its partly built members are destroyed, \cref{ch:b2:cpp-classes-raii}), and no statement has yet
changed \code{*this}: the target keeps its value. If step 1 succeeds, \code{tmp} holds an independent copy of the
value of \code{o}; the swaps give that value to \code{*this} and the old resources to \code{tmp}, whose \omterm{def:b2:cpp-classes-raii:constructor}{destructor}
frees them. When \code{o} is \code{*this}, step 1 copies the target's own value into fresh blocks before anything is
freed, so the result is again that value.
\end{proof}

The self-assignment case is not academic: \code{v[i] = v[j]} with $i = j$ is one, and an assignment written as
``free my block, allocate a new one, copy the source's characters'' reads freed memory in that case
(\cref{exo:b2:cpp-copy-move:2}). The price of copy and swap is a new allocation even when the target's block was
large enough; the standard containers reuse the target's memory when they can, and in exchange do not promise the
strong guarantee for their assignment.

\section{Moving}\label{sec:b2:cpp-copy-move:moving}

A copy is the right meaning of \code{b = a} when \code{a} is used afterwards. When the source is a temporary about
to be destroyed, copying its block and then freeing the original is wasted work: the target could take the block
itself. C++ tells the two cases apart by the kind of expression the source is.

\begin{definition}[Value category, lvalue, rvalue]\label{def:b2:cpp-copy-move:rvalue}
Every expression has a \emph{value category}\index{value category}. An \emph{lvalue}\index{lvalue} names an object
that outlives the expression: a variable \code{a}, \code{*p}, \code{v[i]}. A \emph{prvalue} computes a value that no
name holds yet: \code{42}, \code{x + y}, a call \code{f()} returning by value, \code{Buffer("Quay")}. An
\emph{xvalue} names an object whose resources may be taken, such as \code{std::move(a)}. Prvalues and xvalues together
are the \emph{rvalues}\index{rvalue}.
\end{definition}

\begin{definition}[Rvalue reference]\label{def:b2:cpp-copy-move:rvalue-reference}
An \emph{rvalue reference}\index{rvalue reference} \code{T &&} is a reference that binds only to \omterm{def:b2:cpp-copy-move:rvalue}{rvalues}. When a
function is overloaded on \code{const T &} and \code{T &&}, a call with an \omterm{def:b2:cpp-copy-move:rvalue}{lvalue} picks the first and a call with an
\omterm{def:b2:cpp-copy-move:rvalue}{rvalue} picks the second. \code{std::move(a)} does not move anything: it is a cast of the \omterm{def:b2:cpp-copy-move:rvalue}{lvalue} \code{a} to an xvalue,
which tells the overload resolution that \code{a} may be emptied.
\end{definition}

\begin{definition}[Move constructor, move assignment, moved-from object]\label{def:b2:cpp-copy-move:move}
The \emph{move constructor}\index{move constructor} \code{T(T &&o)} and the \emph{move assignment
operator}\index{move assignment operator} \code{T &operator=(T &&o)} build or assign from an \omterm{def:b2:cpp-copy-move:rvalue}{rvalue} by taking
\code{o}'s resources rather than copying them. Afterwards \code{o} is a \emph{moved-from object}\index{moved-from
object}: it must still be valid, which means it can be destroyed and assigned a new value, but its value is
unspecified unless its class documents it.
\end{definition}

% ledger: b2:cpp11-move
\omterm{def:b2:cpp-copy-move:rvalue-reference}{Rvalue references} and moves arrived in C++11. The \code{Buffer} below has both a copy and a \omterm{def:b2:cpp-copy-move:move}{move constructor}, each
announcing itself; the move takes the pointer and leaves the source with none.

\omcode{code/bachelor-2/02-cpp-copy-move/listings/move_demo.cpp}{14}{45}{A copy, a move and
neither}\label{lst:b2:cpp-copy-move:move-demo}

\omoutput{code/bachelor-2/02-cpp-copy-move/outputs/move_demo.txt}{The move left \code{a}
empty.}\label{out:b2:cpp-copy-move:move-demo}

\begin{omfigure}
\begin{tikzpicture}[font=\footnotesize]
\node[font=\footnotesize\bfseries, anchor=west] at (-0.5,1.9) {before \texttt{Buffer c = std::move(a);}};
\node[anchor=east] at (0,1.0) {\texttt{a}};
\node[cell, minimum width=1.0cm] (an) at (0.55,1.0) {9};
\node[cell, minimum width=1.0cm, right=0mm of an] (ap) {$\bullet$};
\node[draw, fill=omDef!10, minimum width=2.4cm, minimum height=7mm, font=\ttfamily\scriptsize] (blk) at (4.0,1.0) {Mill Lane\textbackslash 0};
\draw[ptr] (ap.center) -- (blk.west);
\begin{scope}[xshift=7.4cm]
\node[font=\footnotesize\bfseries, anchor=west] at (-0.5,1.9) {after};
\node[anchor=east] at (0,1.0) {\texttt{a}};
\node[cell, minimum width=1.0cm] (an2) at (0.55,1.0) {0};
\node[cell, minimum width=1.0cm, right=0mm of an2, font=\ttfamily\scriptsize] (ap2) {null};
\node[anchor=east] at (0,0) {\texttt{c}};
\node[cell, minimum width=1.0cm] (cn) at (0.55,0) {9};
\node[cell, minimum width=1.0cm, right=0mm of cn] (cp) {$\bullet$};
\node[draw, fill=omDef!10, minimum width=2.4cm, minimum height=7mm, font=\ttfamily\scriptsize] (blk2) at (4.0,0) {Mill Lane\textbackslash 0};
\draw[ptr] (cp.center) -- (blk2.west);
\node[font=\scriptsize, text=black!60, anchor=north] at (4.0,-0.45) {same block, new owner};
\end{scope}
\end{tikzpicture}
\par\smallskip{\small A move copies the two members of \code{a} into \code{c} and then empties \code{a}, so that
exactly one object owns the block. Nothing is allocated, and the cost does not depend on the length of the text.}
\end{omfigure}

The third statement printed no \omterm{def:b2:cpp-classes-raii:constructor}{constructor} at all: \code{label} returns the prvalue \code{Buffer(t)}, and C++17
requires that it be built directly in \code{d}.

\begin{definition}[Copy elision]\label{def:b2:cpp-copy-move:elision}
\emph{Copy elision}\index{copy elision} is the construction of an object directly in its final place, which removes a
copy or move. Since C++17 it is guaranteed when an object is initialised from a prvalue of its type: \code{T x =
T(...)}, \code{T x = f()} with \code{f} returning a prvalue, and \code{return T(...);}. When a function returns a
named local object, the compiler may also build it in the caller's object (the \emph{named return value
optimisation}), but is not required to; if it does not, the return moves.
\end{definition}

% ledger: b2:cpp17-elision
This is the first line of \cref{out:b2:cpp-copy-move:hook}: \code{make_readings} returns its local vector \code{v}
by name, the compiler built \code{v} in the caller's \code{r}, and had it not done so the return would have moved the
vector, which copies three pointers and touches no element. The first line counts element copies and moves, so it
would read the same either way: what it shows is that no element was copied, not which of the two happened. Returning a large object by value is therefore cheap in C++, and the usual
style.

\begin{inpractice}[When to write std::move]
Write \code{std::move(x)} only when \code{x} is a named object you will not read again: passing it to a
function that keeps it, or pushing it into a container. Do not write \code{return std::move(local);}: it turns the
returned name into an xvalue, which forbids the elision and forces a move. Do not move from a \code{const} object:
\code{std::move} gives a \code{const T &&}, which binds only to the \omterm{def:b2:cpp-copy-move:copy}{copy constructor}, and the ``move'' silently
copies.
\end{inpractice}

A vector that grows faces the same choice. When it carries its elements to a larger block, moving them is cheaper
than copying them. But if one move throws halfway, some elements are in the new block and some, already moved
from, in the old one, and the strong guarantee of \code{push_back} is lost. The library therefore moves only if the
\omterm{def:b2:cpp-copy-move:move}{move constructor} is declared \code{noexcept}, and copies otherwise. \code{Tracked}, the type of the opening, comes
in two versions that differ only in that word.

\omcode{code/bachelor-2/02-cpp-copy-move/listings/tracked.hpp}{9}{31}{A value that counts its copies and moves, in
two versions}\label{lst:b2:cpp-copy-move:tracked}

\omcode{code/bachelor-2/02-cpp-copy-move/listings/growth.cpp}{6}{20}{Filling a vector from
empty}\label{lst:b2:cpp-copy-move:growth}

\omoutput{code/bachelor-2/02-cpp-copy-move/outputs/growth.txt}{The same pushes, moved or copied at each
regrowth.}\label{out:b2:cpp-copy-move:growth}

% ledger: b2:vector-nothrow-move
Each push moves its temporary into the vector, hence $n$ moves in every line. With $n = 1000$ the capacity went
$1, 2, 4, \dots, 1024$ in 11 regrowths, which carried $1 + 2 + \dots + 512 = 1023$ elements: 1023 more moves with
\code{noexcept}, 1023 copies without it. The count is bounded for any growth factor.

\begin{proposition}[Transfers of a growing array]\label{prop:b2:cpp-copy-move:growth-moves}
An array starts empty and receives $n \geq 2$ elements by appends. Whenever it is full, it moves to a block whose
capacity is at least $r$ times the old one, for a fixed $r > 1$, and carries all its elements over. Then the total
number of elements carried over is less than $\frac{r}{r-1}(n-1)$. For elements whose \omterm{def:b2:cpp-copy-move:move}{move constructor} is
\code{noexcept}, each transfer of a \code{std::vector} is a move and no element is ever copied; otherwise each
transfer is a copy.
\end{proposition}

\begin{proof}
Let $c_1 < c_2 < \dots < c_k$ be the capacities at which the array was full and regrew; regrowth $j$ carries $c_j$
elements. Since $c_{j+1} \geq r c_j$, we have $c_j \leq c_k / r^{k-j}$, and
\[ \sum_{j=1}^{k} c_j \leq c_k \sum_{i=0}^{k-1} r^{-i} < c_k \cdot \frac{1}{1 - 1/r} = \frac{r}{r-1}\, c_k . \]
The last regrowth happened when the array already held $c_k$ elements and received another, so $c_k + 1 \leq n$,
and the sum is less than $\frac{r}{r-1}(n-1)$. If there was no regrowth the sum is $0$. The last sentence is the rule
of the library: it transfers with the \omterm{def:b2:cpp-copy-move:move}{move constructor} exactly when that \omterm{def:b2:cpp-classes-raii:constructor}{constructor} cannot throw, or when there is
no \omterm{def:b2:cpp-copy-move:copy}{copy constructor}.
\end{proof}

For $r = 2$ the bound is $2(n-1)$; the run made 1023 transfers for $n = 1000$, under 1998, and 2047 for $n = 1025$,
under 2048. With $n$ just above a power of two the bound is nearly reached (\cref{exo:b2:cpp-copy-move:11}). Added
to the $n$ insertions themselves, a vector of $n$ elements is filled with fewer than $3n$ element operations, which is
the amortised constant cost recalled at the start, now counted exactly.

\section{The rule of zero and the rule of five}\label{sec:b2:cpp-copy-move:rules}

A class can declare five special members beyond its ordinary \omterm{def:b2:cpp-classes-raii:constructor}{constructors}: the \omterm{def:b2:cpp-classes-raii:constructor}{destructor}, the \omterm{def:b2:cpp-copy-move:copy}{copy constructor},
the copy assignment, the \omterm{def:b2:cpp-copy-move:move}{move constructor} and the move assignment. The compiler writes the missing ones, member by
member, under rules that come down to this: declaring a \omterm{def:b2:cpp-classes-raii:constructor}{destructor} or a copy operation stops the compiler from
writing the moves (copies are used instead), and declaring a move operation deletes the copies.

\begin{definition}[Rule of zero, rule of five]\label{def:b2:cpp-copy-move:rules}
The \emph{rule of zero}\index{rule of zero} says that a class whose members all manage their own resources (strings,
containers, \omterm{def:b2:cpp-copy-move:unique-ptr}{smart pointers}) should declare none of the five special members: the member-wise versions the compiler
writes are correct. The \emph{rule of five}\index{rule of five} says that a class that manages a resource itself, and
so declares one of the five, should declare all five, each written or explicitly \code{= default} or
\code{= delete}.
\end{definition}

\begin{method}[Applying the rule of zero]\label{met:b2:cpp-copy-move:rule-of-zero}
To give a class correct copies and moves:
\begin{enumerate}
\item hold each resource in a member of a type that already owns it: \code{std::string}, \code{std::vector},
\code{std::unique_ptr}, or a small class of your own written once by the \omterm{def:b2:cpp-copy-move:rules}{rule of five};
\item declare none of the five special members in the class itself;
\item check the result: the class copies if every member copies, and moves cheaply if every member moves cheaply.
\end{enumerate}
\end{method}

\omcode{code/bachelor-2/02-cpp-copy-move/listings/rule_of_zero.cpp}{7}{30}{Two classes that declare nothing
special}\label{lst:b2:cpp-copy-move:rule-of-zero}

\omoutput{code/bachelor-2/02-cpp-copy-move/outputs/rule_of_zero.txt}{The copy is deep, the moved-from vector is
empty, the sensor moved.}\label{out:b2:cpp-copy-move:rule-of-zero}

\code{Station} copies and moves correctly with no code at all, because \code{std::string} and \code{std::vector} do.
\code{Sensor} cannot be copied, because one member cannot: a member that is move-only makes its class move-only. The
moved-from \code{b.readings} printed 0, which is what the library's vector leaves behind, though only its validity is
promised. Of the five special members, a class following the \omterm{def:b2:cpp-copy-move:rules}{rule of zero} writes none; \code{Buffer}, which manages
a raw block, is the kind of class that writes them all, once (\cref{exo:b2:cpp-copy-move:3}).

\section{Smart pointers}\label{sec:b2:cpp-copy-move:smart}

In C, a pointer returned by a function says nothing about who must free the block. C++ separates the two meanings
of a pointer: an \emph{owning} pointer, whose holder must free the object, and a plain pointer or reference, which
only refers to an object owned by someone else. Owning pointers are objects of the standard library.

\begin{definition}[Smart pointer, unique pointer, unique ownership]\label{def:b2:cpp-copy-move:unique-ptr}
A \emph{smart pointer}\index{smart pointer} is a class that holds a pointer to a heap object and releases the object
in its \omterm{def:b2:cpp-classes-raii:constructor}{destructor}. The \emph{unique pointer}\index{unique pointer} \code{std::unique_ptr<T>} is the \emph{sole}
owner of its object (\emph{unique ownership}\index{unique ownership}): it cannot be copied, it can be moved, which
transfers the ownership and leaves the source null, and its \omterm{def:b2:cpp-classes-raii:constructor}{destructor} deletes the object if it holds one.
\code{std::make_unique<T>(args)} allocates and constructs the object.
\end{definition}

\omcode{code/bachelor-2/02-cpp-copy-move/listings/unique.cpp}{8}{30}{Photos owned by an
album}\label{lst:b2:cpp-copy-move:unique}

\omoutput{code/bachelor-2/02-cpp-copy-move/outputs/unique.txt}{Each photo is freed when its owner
goes.}\label{out:b2:cpp-copy-move:unique}

A \code{std::unique_ptr} holds one pointer and nothing else, so it costs what a raw pointer costs. The rule that
follows is short: \code{new} and \code{delete} no longer appear in application code. Sometimes, though, an object
has several owners and none of them knows which will finish last: a photo shown in two albums, a cached page held
by several requests.

\begin{definition}[Shared pointer, shared ownership, weak pointer]\label{def:b2:cpp-copy-move:shared-ptr}
A \emph{shared pointer}\index{shared pointer} \code{std::shared_ptr<T>} is one of possibly several owners of an
object (\emph{shared ownership}\index{shared ownership}). The object comes with a \emph{control block} that counts its
shared pointers: copying a \code{shared_ptr} adds one, destroying or resetting one subtracts one, and the owner that
brings the count to zero deletes the object. A \emph{weak pointer}\index{weak pointer} \code{std::weak_ptr<T>}
refers to an object owned by shared pointers without being counted as an owner; \code{lock()} gives a
\code{shared_ptr} to the object if it is still alive, and a null one if not.
\end{definition}

The count makes a \omterm{def:b2:cpp-copy-move:shared-ptr}{shared pointer} cost more than a unique one (a second block, and an update of the count on every
copy), and it has a blind spot. In the program below a photo owns its album and the album owns its cover photo.

\omcode{code/bachelor-2/02-cpp-copy-move/listings/shared_cycle.cpp}{7}{30}{An album and a photo that own each
other}\label{lst:b2:cpp-copy-move:shared-cycle}

\omoutput[1-8,18-22,32-36]{code/bachelor-2/02-cpp-copy-move/outputs/shared_cycle.txt}{LeakSanitizer reports two
blocks nobody can reach.}\label{out:b2:cpp-copy-move:shared-cycle}

Neither \omterm{def:b2:cpp-classes-raii:constructor}{destructor} printed. When \code{main} returned, \code{a} and \code{p} were destroyed and each count fell from
2 to 1, the one held by the other object. Nothing outside the pair points to it any more, but each keeps the other
above zero: the count only sees owners, not whether they are reachable. The fix is to decide that one direction of
the edge observes rather than owns. A photo needs to know its album, but does not keep it alive.

\omcode{code/bachelor-2/02-cpp-copy-move/listings/weak_fix.cpp}{7}{34}{The photo observes its
album}\label{lst:b2:cpp-copy-move:weak-fix}

\omoutput{code/bachelor-2/02-cpp-copy-move/outputs/weak_fix.txt}{Both objects freed; the photo saw its album
go.}\label{out:b2:cpp-copy-move:weak-fix}

\begin{omfigure}
\begin{tikzpicture}[font=\footnotesize, box/.style={draw, rounded corners=2pt, minimum width=2.5cm, minimum height=1.1cm, align=center}]
\node[font=\footnotesize\bfseries, anchor=west] at (-1.3,2.5) {two owning edges};
\node[box, fill=omDef!10] (A) at (0,1.2) {\texttt{Album} Coast\\[1pt] {\scriptsize strong 1, weak 0}};
\node[box, fill=omThm!10] (P) at (3.8,1.2) {\texttt{Photo} pier\\[1pt] {\scriptsize strong 1, weak 0}};
\draw[ptr, bend left=25] (A.north east) to node[above, font=\scriptsize] {cover} (P.north west);
\draw[ptr, bend left=25] (P.south west) to node[below, font=\scriptsize] {album} (A.south east);
\node[font=\scriptsize, text=omThm, align=center] at (1.9,-0.4) {after \texttt{main}: counts stay at 1,\\ nothing is freed};
\begin{scope}[xshift=8.0cm]
\node[font=\footnotesize\bfseries, anchor=west] at (-1.3,2.5) {one weak edge};
\node[box, fill=omDef!10] (A2) at (0,1.2) {\texttt{Album} Coast\\[1pt] {\scriptsize strong 1, weak 1}};
\node[box, fill=omThm!10] (P2) at (3.8,1.2) {\texttt{Photo} pier\\[1pt] {\scriptsize strong 2, weak 0}};
\draw[ptr, bend left=25] (A2.north east) to node[above, font=\scriptsize] {cover} (P2.north west);
\draw[ptr, dashed, black!60, bend left=25] (P2.south west) to node[below, font=\scriptsize] {album (weak)} (A2.south east);
\node[draw, dotted, font=\scriptsize] (va) at (-2.0,1.2) {\texttt{a}};
\draw[ptr, black!60] (va.east) -- (A2.west);
\node[draw, dotted, font=\scriptsize] (vp) at (5.8,1.2) {\texttt{p}};
\draw[ptr, black!60] (vp.west) -- (P2.east);
\node[font=\scriptsize, text=omMeth!80!black, align=center] at (1.9,-0.4) {\texttt{a} goes: album strong 0, freed;\\ then the photo when \texttt{p} goes};
\end{scope}
\end{tikzpicture}
\par\smallskip{\small Each object made by \code{make_shared} carries a control block counting its strong owners and
its weak observers. Left: the cycle of \cref{lst:b2:cpp-copy-move:shared-cycle} once \code{a} and \code{p} are gone.
Right: inside the block of \cref{lst:b2:cpp-copy-move:weak-fix}, while \code{a} and \code{p} still exist; the dashed
edge does not count, so the album dies with \code{a}.}
\end{omfigure}

\begin{inpractice}[Choosing an owner]
Most objects should have exactly one owner, and \code{std::unique_ptr} or a plain member says so. Reach for
\code{std::shared_ptr} only when the lifetime is genuinely shared, and then draw the ownership graph: every cycle
in it must contain a weak edge. Pointers back to a parent, to an observer, or to a cache entry are the usual
candidates. A child owned by its parent can even point back with a raw pointer, since the parent always outlives it
(\cref{exo:b2:cpp-copy-move:9}). Book 5 returns to reference counting as a technique for managing memory, and to its
alternatives.
\end{inpractice}

Moving costs a few pointer copies whatever the size of the object; copying costs time proportional to the data. The
program below measures both for vectors of \code{double} of growing length.

\omcode{code/bachelor-2/02-cpp-copy-move/listings/copy_vs_move.cpp}{16}{28}{Copy and move of a vector, timed
(\code{ns} returns the mean time of \code{reps} calls of its argument)}\label{lst:b2:cpp-copy-move:copy-vs-move}

\omoutput{code/bachelor-2/02-cpp-copy-move/outputs/copy_vs_move.txt}{One run on the author's laptop; the digits
change from run to run.}\label{out:b2:cpp-copy-move:copy-vs-move}

The copy grows with $n$: from a few hundredths to about a tenth of a nanosecond per element while the vectors fit in
the processor's caches, and about half a nanosecond once they do not (\cref{ch:b2:caches}). The move stays below two
nanoseconds whatever the length, because it copies three pointers and empties the source.

\section{Exercises}\label{sec:b2:cpp-copy-move:exercises}

\begin{exercise}[$\star$]\label{exo:b2:cpp-copy-move:1}
Using \code{Tracked} from \cref{lst:b2:cpp-copy-move:tracked}, say for each of the lines A to F whether it copies,
moves, or neither, and give the two numbers the program prints.
\omcode{code/bachelor-2/02-cpp-copy-move/listings/exo_count.cpp}{7}{22}{Copies and moves, line by line}
\end{exercise}

\begin{exercise}[$\star$]\label{exo:b2:cpp-copy-move:2}
The class \code{Text} below has a copy assignment written the obvious way. The program assigns an object to
itself. Run under AddressSanitizer, it stops at once. Explain the report, then fix the assignment twice: with a test
for self-assignment, and with \cref{met:b2:cpp-copy-move:copy-and-swap}. Which fix also gives the strong guarantee?
\omcode{code/bachelor-2/02-cpp-copy-move/listings/exo_selfassign.cpp}{5}{32}{An assignment to oneself}
\omoutput[1-8]{code/bachelor-2/02-cpp-copy-move/outputs/run_selfassign.txt}{The beginning of the report.}
\end{exercise}

\begin{exercise}[$\star$]\label{exo:b2:cpp-copy-move:3}
Complete the \code{Buffer} of \cref{lst:b2:cpp-copy-move:buffer} by the \omterm{def:b2:cpp-copy-move:rules}{rule of five}: add a \code{noexcept} \omterm{def:b2:cpp-copy-move:move}{move
constructor} that leaves the source empty, and replace the copy assignment by a single assignment
\code{operator=(Buffer o)} that takes its parameter by value. Explain why this one function serves as both copy and
move assignment.
\end{exercise}

\begin{exercise}[$\star$]\label{exo:b2:cpp-copy-move:4}
A vector doubles its capacity from 1 whenever it is full. Count exactly the elements carried over while it
receives $n = 2^m$ elements, then $n = 2^m + 1$. Check both against \cref{prop:b2:cpp-copy-move:growth-moves} and
against \cref{out:b2:cpp-copy-move:growth}.
\end{exercise}

\begin{exercise}[$\star\star$]\label{exo:b2:cpp-copy-move:5}
For each of the eight numbered statements, say which special member of \code{W} runs, or that none does, and
which rule of the chapter decides it. One of the eight depends on the compiler: which one, and what may another
compiler print?
\omcode{code/bachelor-2/02-cpp-copy-move/listings/exo_which.cpp}{4}{28}{Eight statements}
\end{exercise}

\begin{exercise}[$\star\star$]\label{exo:b2:cpp-copy-move:6}
The program below should keep, for each day, the day's list of jobs. Explain the output. Is it undefined behaviour?
Fix the loop so that each day starts from the same three jobs, and say which line of the fix copies.
\omcode{code/bachelor-2/02-cpp-copy-move/listings/exo_after_move.cpp}{6}{18}{A job list kept each day}
\omoutput{code/bachelor-2/02-cpp-copy-move/outputs/exo_after_move.txt}{The output.}
\end{exercise}

\begin{exercise}[$\star\star$]\label{exo:b2:cpp-copy-move:7}
Write a class template \code{UniquePtr<T>} with the behaviour of \cref{def:b2:cpp-copy-move:unique-ptr}: a
\omterm{def:b2:cpp-classes-raii:constructor}{constructor} from \code{T *}, no copies, a \code{noexcept} \omterm{def:b2:cpp-copy-move:move}{move constructor} and move assignment (a self-move must not
delete the object), \code{*}, \code{->}, \code{get}, \code{release}, \code{reset} and a test for null. Test that a
vector of a hundred \code{UniquePtr} frees each object exactly once.
\end{exercise}

\begin{exercise}[$\star\star$]\label{exo:b2:cpp-copy-move:8}
Prove \cref{prop:b2:cpp-copy-move:copy-and-swap} again in the form of \cref{exo:b2:cpp-copy-move:3}, where the
parameter is taken by value. Where does the copy that may throw now happen, and why is the target still unchanged
if it throws? Does the move assignment that this function also provides give the strong guarantee?
\end{exercise}

\begin{exercise}[$\star\star$]\label{exo:b2:cpp-copy-move:9}
Represent a tree of named places (town, quarters, buildings) with nodes that own their children through
\code{std::vector<std::unique_ptr<Node>>} and point to their parent with a raw \code{Node *}. Write \code{add},
\code{path} (\code{town/west/mill}) and \code{move_child(i, to)}, which moves a whole subtree under another node.
Why is a raw parent pointer safe here, and why would a \code{std::shared_ptr} to the parent leak? Does the moved
subtree change address?
\end{exercise}

\begin{exercise}[$\star\star\star$]\label{exo:b2:cpp-copy-move:10}
Write a string class \code{SmallString} that keeps strings of at most 15 characters inside the object, in a member
array, and longer ones in a heap block (a \emph{small-buffer optimisation}). Give it the five special members. What
must the \omterm{def:b2:cpp-copy-move:move}{move constructor} do when the source is small, and why is it then not cheaper than a copy? Test copies and
moves between every combination of small and long strings, self-assignment, and growth of a vector of them.
\end{exercise}

\begin{exercise}[$\star\star\star$]\label{exo:b2:cpp-copy-move:11}
(a) Show that the bound of \cref{prop:b2:cpp-copy-move:growth-moves} cannot be improved to $\frac{r}{r-1}(n-1)
\cdot \alpha$ for any $\alpha < 1$: exhibit, for $r = 2$, values of $n$ for which the number of transfers divided by
$2(n-1)$ tends to $1$. (b) What does the proof give for $r = 1.5$, and how many transfers does a growth from 1 by
the factor 1.5 rounded down, but by at least one place (capacities $1, 2, 3, 4, 6, 9, \dots$), make up to $n = 100$? (c) An array that grows by a
fixed number $d$ of places instead of a factor makes how many transfers for $n$ elements? Conclude about the
amortised cost.
\end{exercise}

\begin{exercise}[$\star\star\star$]\label{exo:b2:cpp-copy-move:12}
A station keeps the displays that show its readings, and each display keeps its station. Explain the output, draw
the ownership graph with the counts at the end of \code{main}, and fix the program with one \code{std::weak_ptr}.
Which edge should it be, and why not the other? Prove that, in your fixed version, every object is freed when the
last outside \code{shared_ptr} to the station goes, for any number of displays.
\omcode{code/bachelor-2/02-cpp-copy-move/listings/exo_observers.cpp}{6}{37}{A station and its displays}
\omoutput[1-5,63-64]{code/bachelor-2/02-cpp-copy-move/outputs/exo_observers.txt}{The start and the end of the
report.}
\end{exercise}

\section{Problem: The photo library that copied itself}\label{sec:b2:cpp-copy-move:problem}

\begin{problem}[{Weekend problem --- The photo library that copied itself}]\label{pb:b2:cpp-copy-move:1}
A photo library edits images of $4000 \times 3000$ pixels, one byte per pixel, held by value in an \code{Image} class
that counts its copies and moves. It runs the same two-step edit in two styles and then files three images in an
album.

\omcodes{code/bachelor-2/02-cpp-copy-move/listings/pipeline.cpp}{9-22,24-29,31-49}{An editing
pipeline}\label{lst:b2:cpp-copy-move:pipeline}

\needspace{4\baselineskip}\medskip\noindent\textbf{Part I --- Counting copies.}\par\nopagebreak
\begin{enumerate}[beginpenalty=10000]
\item How many bytes does one copy of an \code{Image} allocate and copy?
\item In version 1, which statements copy an \code{Image} when \code{a} is computed? Predict the count of copies and
moves for the first report.
\item Why does \code{return out;} in \code{brighten1} not count as a move in this run? What would it count with a
compiler that does not apply the named return value optimisation?
\item In version 2, where does the one copy of the second report happen, and why must it happen if \code{photo} is
used again?
\item Where do the two moves of the second report happen? Why is \code{return im;} a move and not elided?
\item Why does the result of \code{brighten2} reach the parameter of \code{darken2} with neither a copy nor a move?
\end{enumerate}

\needspace{4\baselineskip}\medskip\noindent\textbf{Part II --- Adding moves.}\par\nopagebreak
\begin{enumerate}[resume, beginpenalty=10000]
\item Predict the third report (with \code{std::move(photo)}). What may the program do with \code{photo} afterwards?
\item Predict the fourth report: the three \code{push_back} calls, with the vector growing from capacity 0 as in
\cref{out:b2:cpp-copy-move:growth}.
\item How many bytes are copied in the whole run? Compare with a version where \code{Image}'s \omterm{def:b2:cpp-copy-move:move}{move constructor} lacks
\code{noexcept}.
\item Rewrite the album part with \code{album.reserve(3)}; what does the fourth report become?
\item \code{Image} declares its five special members only to count. Which of them would the \omterm{def:b2:cpp-copy-move:rules}{rule of zero} let it
drop, and what would the class then do on a copy and a move?
\item State a rule for choosing between \code{const Image &} and \code{Image} as the parameter type of an edit
function, from the counts above.
\end{enumerate}

\needspace{4\baselineskip}\medskip\noindent\textbf{Part III --- Albums and photos.}\par\nopagebreak
\begin{enumerate}[resume, beginpenalty=10000]
\item The library now lets one photo belong to several albums, and an album keeps its photos with
\code{std::vector<std::shared_ptr<Photo>>}. Why is \code{std::unique_ptr} no longer possible?
\item Each photo keeps the list of albums that contain it, as \code{std::vector<std::shared_ptr<Album>>}. Draw the
ownership graph for two albums sharing one photo, with every strong count, once the library's own pointers to the
albums are gone.
\item Which objects leak? How many bytes, for photos of the size above?
\item Fix it with \omterm{def:b2:cpp-copy-move:shared-ptr}{weak pointers}. Which edges become weak, and what are the counts afterwards?
\item When an album is deleted, a photo's list keeps an expired \code{weak_ptr}. Write the loop that removes the
expired entries, using \code{expired()} or \code{lock()}.
\item A photo shown on screen must not disappear while it is being drawn, even if its last album is deleted at that
moment by another part of the program. Which call gives the drawing code that guarantee?
\end{enumerate}

\needspace{4\baselineskip}\medskip\noindent\textbf{Part IV --- What a copy costs.}\par\nopagebreak
\begin{enumerate}[resume, beginpenalty=10000]
\item In \cref{out:b2:cpp-copy-move:copy-vs-move}, estimate the time per element of a copy for $n = 10^3$ and for
$n = 10^7$. Why do they differ?
\item Estimate the time of one copy of an \code{Image} from those figures (its elements are bytes, not
\code{double}s: state your assumption).
\item The library processes 50\,000 photos a night with version 1 and then files them. Estimate the time spent in
copies, and the time saved by version 2 with \code{std::move}.
\item Why does the move column not grow with $n$? Name the members a move of \code{std::vector} touches.
\item A vector of $n$ images grows by doubling, without \code{reserve}. With \cref{prop:b2:cpp-copy-move:growth-moves},
bound the bytes copied if \code{Image}'s moves were not \code{noexcept}, for $n = 1000$.
\item State the result of the problem: the number of transfers of a growing array, and what each costs with and
without \code{noexcept} moves.
\end{enumerate}
\end{problem}
